% Budget: ~1 column  |  10 pts.  Slides p.12-13.
% Sources: data/README.md (P1), research_story/README.md (I1), model notebooks (P3).
\section{Experimental Setup}
\label{sec:setup}

\subsection{Validation Protocol}
The 292 test sales are held out from the start and used once, after the features, model, hyperparameters and threshold are frozen.
The 1{,}166 training sales form 5 fixed folds (233/234/234/232/233 rows) shared by every method; each row is a distinct sale predicted at the property level, not a time series, so we use plain K-fold rather than a grouped or temporal split.
Five folds keep about 930 rows for fitting and over 230 per validation fold; we report mean $\pm$ SD.

\subsection{Metrics and Baselines}
\textbf{Metrics.} MAE in USD is the primary metric, because it is the typical size of a pricing mistake a seller or buyer faces.
RMSE in USD emphasizes large misses, which are costly in a negotiation.
$R^2$ shows the share of price variance explained.

\textbf{Decision threshold.} A model is usable if its MAE is at most 15{,}000~USD, fixed before testing, for two reasons.
For users, it is about 9\% of the median training price, accurate enough for a first asking-price estimate but not a substitute for a formal appraisal.
For generalization, it is the final model's CV MAE plus two fold SDs ($13{,}366 + 2 \times 809 = 14{,}985$~USD), rounded up, so a higher test MAE would signal worse performance than normal fold variation.
All three baselines exceed it in CV, so it separates usable from unusable models.

\textbf{Baselines.} (i)~\emph{Trivial:} the training-set mean price.
(ii)~\emph{Simple:} OLS on living area alone, the most intuitive single feature.
(iii)~\emph{Strong:} a Random Forest (300 trees) on the raw features and raw price, tuned over maximum depth $\{\text{none}, 10\}$, minimum leaf size $\{1, 3\}$ and feature fraction $\{0.7, 1.0\}$.
We also tried the log target for the forest; it was worse (CV MAE 16{,}800 vs.\ 16{,}554~USD), so we keep the raw price.
